% Legacy analyses of the earlier version of the study, moved out of the supplement on 2026-10-05 (verbatim; cross-references to the supplement may no longer resolve).
% Compile only as part of a document that defines the supplement's macros.

\subsection{Additional checkpoints and cameras}
\label{app:additional}

The remaining analyses extend the coverage of the main result using transfer
checkpoints and public models. They are retained as secondary evidence and
do not supply a prerequisite result from another paper. Model provenance and
the Human-M3 source-frame overlap are disclosed in \ref{app:shared}.

\subsection{The original checkpoint on Q1\_top: a cross-view observation}
\label{app:q1_crossview}

Q1\_top is an unused camera view for the original CourtDyn-native checkpoint,
but its opening event also appears in that checkpoint's Q1\_side training clip
(event-level split audit of the main text). It is therefore a cross-view observation and not
an independent event-level test, and no claim in the main text rests on it. We
report it here because it was collected under exactly the legacy protocol of
the main text and because withholding a measured cell would be its
own distortion. The matched adapters of
Section~\ref{sec:matched_label_results} evaluate the same clip without this
overlap, and those results---not the ones below---carry the Q1 evidence.

Under v3 metre scoring the speed/path T-MRA scores are $69.3$/$52.9$ against
test-distribution constants of $43.0$/$35.0$; pixel-unit scores are
$16.6$/$12.1$ against $43.8$/$35.1$, with paired pixel/full ratios of
$1.08$/$0.62$ and pixel-unit correlations of $0.68$/$0.60$. Under the
first-frame repetition control the metre correlations change from
$0.71$ to $0.02$ and from $0.68$ to $-0.04$ ($\Delta\rho=+0.69$/$+0.72$),
with distinct predictions falling from $11$ and $13$ to $3$.

The direction of every one of these effects matches the event-disjoint Q2 cell
of the main text, which is the reason we treat the overlap as
affecting the strength of the evidence rather than its sign. Combining the two
clips would let a contaminated cell widen a reported range, so ranges in the
main text are computed from Q2 alone. Figure~\ref{fig:nativecore_both} plots
both clips together for completeness, with the overlapped pair shaded.

\begin{figure}[t]
\centering
\includegraphics[width=\linewidth]{fig_native_core.pdf}
\caption{Both clips of the original checkpoint under the legacy prompts:
  (a) advantage over the test-distribution constant, (b) median paired
  ratios. Shading marks the Q1 cells, whose opening event appears in
  training through another camera. The main text uses the unshaded cells
  only.}
\label{fig:nativecore_both}
\end{figure}

\subsection{Which ruler scores it, and why the answer matters}
\label{sec:whichruler}

The natural next question---does the model have a sense of absolute
scale?---turns out to depend on a choice the benchmark designer makes, and the
dependence is large enough to be a finding in its own right.

\begin{figure}[t]
\centering
\includegraphics[width=\linewidth]{fig2_constant_vs_gt.pdf}
\caption{Model T-MRA minus that clip's constant-median baseline, plotted against
  the baseline itself, under both ground truths ($56$ cells each: 7 clips
  $\times$ 4 models $\times$ 2 families). Points above zero beat a predictor
  with no image input but access to the test answer distribution. The horizontal
  coordinate is also subtracted in the vertical coordinate, so any downward
  trend is not by itself evidence of a causal effect of distribution width. Plotted from recorded outputs and stated references using AI-assisted Python code (see the AI-use declaration of the main text).}
\label{fig:constant}
\end{figure}

Recomputing every cell under both ground truths with the same scoring function
(Figure~\ref{fig:constant}),
$8$ of $56$ cells exceed their own clip's constant baseline under v1 and $9$ of
$56$ under v3; the base model exceeds it in none. That agreement is the robust
part. The disagreement is not: individual cells move substantially and some
change sign. The source-SFT path cell on Q1\_top is $10.9$ points \emph{below}
the constant baseline under v1 and $13.9$ points \emph{above} it under v3.

The relationship between rank agreement and tolerance score therefore depends
on the scoring convention. Under v3, some high-correlation cells also exceed
the constant baseline. The supported methodological observation is that a
box-height heuristic and a court-grounded reference can reverse evaluations
of the same predictions. A prompt stating a typical height does not uniquely
prescribe the v1 calculation. We therefore recommend reporting both the
prompted scale cue and the scoring convention rather than attributing the
entire discrepancy to instruction following.

Two results about scale are independent of this choice, because both are
within-arm comparisons.

\emph{Deleting the height sentence.} We delete ``assume a typical player is
1.93\,m tall'' from the prompt entirely. Across all three overhead clips and
both adapters (12 cells), the predicted median moves by a factor of $0.97$ to
$1.10$, and in seven of the twelve cells it does not move at all: $0.40 \to
0.40$, $1.00 \to 1.00$, $1.90 \to 1.90$, $4.10 \to 4.10$. This is limited sensitivity of the output median to deleting this cue;
it does not establish whether the cue was internally represented.

One cell deserves recording rather than averaging away. On Q1\_top the
source-SFT speed correlation falls from $\rho = 0.59$ to $0.43$, the only
sizeable \rhosp{} movement in the twelve, while its predicted median is
unchanged at $0.40$. Thus the rank statistic and the output median can respond
differently to deleting the sentence.

\emph{Post-hoc no-prior decomposition.} For source-SFT and GRPO-v3 on
Q1/Q2 speed and path, all eight cells have 140 common finite predictions in
the full, no-prior and pixel arms, with positive ratio denominators. The
no-prior prompts differ from full only by deleting the height sentence;
images and inference configurations match. Median paired no-prior/full
ratios range from $0.983$ to $1.000$, but deleting the sentence changes v3
T-MRA by $-3.50$ to $+4.93$ points. Median paired pixel/no-prior ratios
remain $1.00$--$1.79$, far below the nominal rendered-coordinate factor
of about 27. Thus deleting the height sentence alone does not explain the
weak pixel response in these two auxiliary models. The pixel arm still
changes quantity wording and leaves coordinates unspecified. The legacy
no-prior outputs used here cover source-SFT and GRPO-v3; the new native
controls are reported separately in the main text.

\emph{A counterfactual height cue.} Replacing the stated height of 1.93 m by 1.45 m and rescaling the v1
speed/path reference by $0.751$ tests responsiveness to that heuristic.
A predictor explicitly implementing v1 would scale its numeric output likewise. The ratios of prediction medians are $1.692$ for
base, $1.000$ for source-SFT and $1.100$ for GRPO, and source-SFT's score
\emph{rises} ($+13.9$ speed, $+7.7$ path).

\subsection{Scale and unit comparisons across checkpoints}
\label{sec:scale_transfer}

We test an explicit scale cue and a requested-unit change on two overhead
clips. Both arms use the same source frames and item identities as full.
The scale cue and pixel ground truth use the $960\times540$ rendered canvas;
this convention was not specified explicitly in the original model prompts.

\begin{figure}[t]
\centering
\includegraphics[width=\linewidth]{fig5_ruler_and_units.pdf}
\caption{Median paired prediction ratios for the four adapters.
  (a) A clip-level scale cue yields ratios of $0.67$--$1.00$.
  (b) Pixel-unit prompts yield $0.62$--$2.00$, relative to approximately
  $K\approx27$ under the rendered-image convention. Panel (b) uses a
  logarithmic axis. These ratios do not identify an internal mechanism. Plotted from recorded outputs and stated references using AI-assisted Python code (see the AI-use declaration of the main text).}
\label{fig:ruler}
\end{figure}

\paragraph{The ruler arm} We replace the height sentence with
``In these frames, one meter on the floor spans about 27 pixels.'' The underlying
clip-level $K$ is the median rendered-pixel path length divided by homography
path length over eligible path items; it is not a per-item box-height cue.
The numerical factors are $27.061$ and $27.498$ pixels/m, both rounded to 27
in the prompts. The builder requires an interquartile range divided by the
median of at most $5\%$. Full and ruler predictions are both scored against
v3. This controls the scoring convention but not ambiguity about model-side
image resizing.

For the three transfer checkpoints, $8$ of $12$ cells have a median paired
prediction ratio of $1.00$; the remaining four give $0.79$--$0.96$. $\Delta$T-MRA scatters between $-8.1$ and
$+13.0$ with no systematic improvement. Supplying this cue does not systematically improve the score. This does not
separate coordinate interpretation, instruction following and arithmetic.
The base model is the apparent exception (ratios $1.88$ and $12.00$), but it
simultaneously drives \rhosp{} from $+0.11$ to $-0.19$; its changing magnitude does not accompany improved rank agreement.
Figure~\ref{fig:ruler} presents the four adapters, including CourtDyn-native.

\paragraph{The unit-swap arm} We re-ask the identical items, on identical
frames, with the answer unit and the ground truth expressed in rendered pixels.
The tolerance is scaled by $K$ so the two arms remain comparable.
Table~\ref{tab:unitswap} reports the outcome.

\begin{table}[t]
\centering
\caption{Unit swap, per clip (Q1 / Q2), $K = 27.06$ / $27.50$ rendered pixels
  per metre. Substantial rank agreement remains for adapted models; $R$ denotes the
  median paired ratio in the main text.}
\label{tab:unitswap}
\small
\begin{tabular}{llll}
\toprule
Model & Family & Pixel-unit \rhosp{} & $R$ (pixel/full) \\
\midrule
source-SFT        & speed & $0.52$ / $0.58$ & $2.00$ / $1.50$ \\
source-SFT        & path  & $0.77$ / $0.66$ & $1.00$ / $1.00$ \\
GRPO-v3           & speed & $0.64$ / $0.57$ & $1.30$ / $1.36$ \\
GRPO-v3           & path  & $0.75$ / $0.68$ & $1.03$ / $1.00$ \\
target-SFT-token  & speed & $0.52$ / $0.64$ & $1.30$ / $1.10$ \\
target-SFT-token  & path  & $0.77$ / $0.71$ & $1.27$ / $1.25$ \\
\midrule
Qwen2.5-VL-7B     & speed & $0.13$ / $0.00$ & $0.87$ / $0.87$ \\
Qwen2.5-VL-7B     & path  & $-0.01$ / $-0.14$ & $1.00$ / $1.00$ \\
\midrule
CourtDyn-native   & speed & $0.68$ / $0.70$ & $1.08$ / $1.11$ \\
CourtDyn-native   & path  & $0.60$ / $0.73$ & $0.62$ / $0.62$ \\
\midrule
\multicolumn{2}{l}{\emph{nominal rendered-image reference $K$}} & --- & $27.1$ / $27.5$ \\
\bottomrule
\end{tabular}
\end{table}

The transfer checkpoints retain positive rank agreement ($0.52$--$0.77$),
compared with $0.34$--$0.80$ under v3 metre scoring. These ranges do not prove
unchanged correlations. Median paired prediction ratios of $1.00$--$2.00$
remain far below approximately $27$ under the rendered-image convention.
Pixel-unit ground-truth medians are $29.6$--$88.5$; the corresponding model
output medians are $0.6$--$4.9$.

This is evidence for weak responsiveness to the requested unit while retaining
ordering. The original prompts do not state which image coordinate system to
use after preprocessing. Thus $K$ defines our evaluation convention rather
than a uniquely communicated model-side conversion. The result does not
identify an absent internal magnitude representation.

\paragraph{Public-model comparison}
We also ran public Qwen2.5-VL-7B through both arms. The handed ruler
gives median paired prediction ratios of $0.87$ in all four ruler cells
and $0.87$--$1.00$ in the unit swap, compared with approximately $27$ under
the rendered-image convention. This particular public model does not provide a successful unit-conversion
comparison; its task-level degeneracy precludes conclusions about stronger models.

It does not, however, independently replicate the behavioural conjunction, and we will not
present it as though it did. The 7B is \emph{degenerate} on these arms: it emits
essentially a single number (median $1.5$ in metres, $1.3$ once a ruler is
handed over), the distinct-prediction count falls to one or two, and Spearman's
\rhosp{} is undefined in three of the four ruler cells for want of variation.
Its pixel-unit \rhosp{} is $-0.14$ to $0.13$---indistinguishable from zero,
unlike the $0.52$--$0.77$ of the adapters. This weak task performance does not independently reproduce the conjunction
of retained ordering and limited unit response.

The supported conclusion is restricted to these prompts, camera clips and
adapters: positive rank agreement coexists with unit responses much smaller
than the rendered-pixel reference factor. Output priors and instruction or
coordinate interpretation remain alternative explanations.

\paragraph{Token-level supervision} The target-SFT-token arm has zoom ratios
of $1.16$--$1.81$ and first-frame sensitivity of similar magnitude to the other
transfer checkpoints. This arm supplies no evidence that its particular static
supervision recipe resolves unit responsiveness; it does not test every form
of token-level or multi-unit supervision.

\subsection{Pixels or metres: the zoom intervention}
\label{sec:zoom}

\begin{figure}[t]
\centering
\includegraphics[width=\linewidth]{fig4_zoom_ratio.pdf}
\caption{Median $\mathrm{pred(zoom)}/\mathrm{pred(full)}$ per model and family,
  one point per clip. Ratios are reported per clip against the planned reference criteria.
  The 7B group is shaded: its ratio of $1.00$ is degenerate. Plotted from recorded outputs and stated references using AI-assisted Python code (see the AI-use declaration of the main text).}
\label{fig:zoom}
\end{figure}

Under P3, doubling pixel displacement while holding the world quantity fixed
gives median paired prediction ratios of $1.42$ and $1.05$ (source-SFT speed and path),
$1.21$ and $1.00$ (GRPO), and $1.00$ and $0.87$ (Qwen2.5-VL-7B), summarised in
Figure~\ref{fig:zoom}. The fraction of items whose ratio exceeds $1.5$ is
$34\%$, $21\%$, $33\%$, $10\%$, $15\%$ and $5\%$ respectively.

The planned binary criteria did not yield a decisive classification, so we
report the distribution. These results do not by themselves exclude either
mechanistic endpoint: mixtures, output quantisation and degenerate constant
answers can yield intermediate or near-unit ratios. Concentration on a few
numeric values is consistent with an answer-prior account, without identifying
it uniquely.

The 7B's ratio of $1.00$ must not be read as ruler use: it emits a constant
$1.5$ in both arms, so the ratio is degenerate. The criterion is only meaningful
where the predictions vary.

\paragraph{An unplanned result: cropping restores the side views}
The zoom arm also raises \rhosp{} itself, and far more on side views than
overhead: source-SFT speed $0.42 \to 0.57$, path $0.54 \to 0.70$; GRPO
$0.47 \to 0.58$ and $0.53 \to 0.69$; the 7B path arm $0.01 \to 0.35$. The median
$\Delta\rho$ is $+0.06$ overhead and $+0.36$ on side views, where the same
clip's \rhosp{} rises from $0.0$--$0.2$ to $0.43$--$0.50$---the overhead range.

These results show that side-view rank agreement can improve after the
combined crop-and-magnification intervention. They do not isolate projective
geometry from apparent size and surrounding content. Centring the crop on
the trajectory both magnifies the target and deletes everything outside the
crop, so we ran a third arm to compare these two components.

\paragraph{Decomposing the crop}
The \emph{declutter} arm applies the identical crop rectangle (the same function
produces it, so the retained pixels and the dropped items match the zoom arm
item for item) but renders it at the \emph{full} arm's scale, $480 \times 270$,
centred on a $960 \times 540$ black canvas. Apparent size is then exactly that
of the full arm and the retained content exactly that of the zoom arm, so with
all three arms scored on their common items,
$\rho(\text{declutter}) - \rho(\text{full})$ compares the content-and-layout change
and $\rho(\text{zoom}) - \rho(\text{declutter})$ compares magnification conditional on that crop.

On the side clip the split is lopsided: the content term has median $+0.260$
(range $+0.16$ to $+0.34$ over the four model $\times$ family cells) and the
size term median $+0.081$ ($+0.06$ to $+0.17$). The content contrast accounts descriptively for roughly $70\%$ of the
recovery along this sequence of interventions. It is not an independently
identified causal fraction: the black canvas also changes layout and centring. On the three overhead clips both terms are near zero (medians
$+0.071$ and $+0.005$), consistent with the small overhead effect above; in
five of twelve overhead cells the size term is in fact \emph{negative}, i.e.\
magnification slightly hurts.

This comparison suggests a contribution beyond magnification alone;
it does not establish independent effects of resolution and perspective geometry. One boundary must be kept in the wording. The
black border removes \emph{everything} outside the crop---distracting players,
but also the court lines and the background, which are exactly the perspective
cues. The content term is therefore the effect of \emph{surrounding content},
not of \emph{other players}; isolating the players alone would need a fourth arm
that paints them out while keeping the court, which we did not run.

Note also that this does not conflict with the item-quality audit of
the main text: removing low-quality \emph{items} barely moves
\rhosp{}, whereas changing the \emph{image} moves it a great deal. The two
operate on different objects.

\subsection{Paired families}
\label{sec:paired}

On the vision-side minimal pairs (chance $25$), the protocol behaves as
designed. Under a blank image the paired score is exactly $0.0$, because a constant answer necessarily scores zero when the pair is scored
jointly. This conclusion follows from the pair construction and the present
blank-image evaluations.

The capability itself, however, is not established. Paired accuracy on
accel\_vis is $3.4$ (base), $22.4$ (source-SFT), $25.9$ (GRPO) and $32.8$
(target-SFT-token); magnifying by $4.5\times$ does not rescue it
($19.0 / 15.5 / 17.2 / 24.1$). On faster, only Q2\_top rises meaningfully
(GRPO $41.4$, 7B $32.9$, SFT $30.0$); Q1\_top is at most $17.2$ everywhere.
These results do not establish a dynamics-comparison ability: most cells sit
at or below chance.

\subsection{Reproduction arms}
\label{sec:repro}

To assess the scope beyond our adapters, two public models were run
through the same protocols. Qwen2.5-VL-7B reaches $\rho \le 0.25$ on every clip,
shows only small first-frame differences ($\Delta\rho = 0.06$--$0.28$) and a
zoom ratio locked at $1.00$. SmolVLM2-2.2B scores \emph{higher} on a blank image
than on the real one (speed T-MRA $76.9$ vs $31.4$) and $0.0$ paired under a
blank.

Both show weak task performance compatible with an answer-prior account.
Neither establishes the combination of positive rank agreement and weak unit
responsiveness observed for the task-trained model. These results do not
show that the public models use no multi-frame content. They are zero-shot
comparisons; the trained replications on SmolVLM2-2.2B, Qwen2.5-VL-3B and
InternVL3-2B, with decision rules recorded before training, are reported in the
main text.

\subsection{A second venue and an inconclusive camera comparison}
\label{sec:secondvenue}

Everything above comes from one basketball venue (the soccer evaluation of the
main text is a second sport). On Human-M3~\cite{humanm3} basketball
sequences, we evaluated source-SFT, GRPO and a public comparison model.
This is an auxiliary transfer
analysis, not an external replication of CourtDyn-native. Ground truth is the
dataset's own 3D world coordinates---no height ruler, no homography, and hence
none of the factor-of-two ambiguity of \ref{sec:whichruler}. Crucially,
the same world motion is captured simultaneously by several cameras with
different geometry, so viewpoint can be varied on \emph{identical} items with
byte-identical ground truth.

\begin{figure}[t]
\centering
\includegraphics[width=\linewidth]{fig3_multiview.pdf}
\caption{Human-M3: coupling between pixel and world displacement per camera,
  with model rank correlations overlaid. Plotted from recorded outputs and stated references using AI-assisted Python code (see the AI-use declaration of the main text).}
\label{fig:multiview}
\end{figure}

\paragraph{Observed behaviour} Figure~\ref{fig:multiview} overlays the model
correlations on the per-camera pixel--world coupling. In this auxiliary
evaluation, across 28 model $\times$ family $\times$ camera rows, $26$ fall below
their own camera's constant baseline (deficits $-1.9$ to $-63.5$), and the two positive values are $+2.9$ and $+0.6$. We have not established
their statistical significance. The predominance of deficits also occurs
outside the CourtDyn homography setting.

The transfer checkpoints also show a first-frame effect. On the camera with the
highest full-frame \rhosp{} (basketball2, camera~0), removing motion lowers
\rhosp{} in all four cells: $0.48 \to 0.28$, $0.50 \to 0.13$, $0.44 \to 0.18$
and $0.46 \to 0.17$, i.e.\ $\Delta\rho$ of $0.20$, $0.37$, $0.26$ and $0.29$.
The magnitude, however, is about half of what the overhead CourtDyn clips give
($0.46$--$0.93$), the internally planned threshold of $0.3$ is cleared by one cell
of four, and the control's \rhosp{} does not reach zero ($0.13$--$0.28$ versus
$0.01$ to $-0.17$ overhead). We therefore write that multi-frame use replicates
\emph{in direction at roughly half the size}, and not that it replicates fully.

\paragraph{Camera criterion} An internal plan recorded before running specified the
prediction that model \rhosp{} should track the per-camera coupling between
pixel and world displacement, with the alternative that flat model \rhosp{}
across cameras would challenge the image-plane interpretation in that plan.
Neither outcome was obtained. The rank correlations between model \rhosp{} and that coupling are
$+0.63$, $-0.40$, $+0.40$ and $-0.06$ across the four model--family
combinations: two are negative. The camera with the weakest coupling ($0.13$)
is indeed the one where both models' speed correlation collapses to $-0.00$ and
$0.03$, which points the right way; but the camera with the strongest coupling
($0.77$) shows model \rhosp{} of only $0.17$--$0.24$, the lowest in the table,
which points the wrong way. Nor is model \rhosp{} flat (speed ranges $0.00$ to
$0.34$ within one sequence).

This is a rank comparison over only four cameras. Its mixed signs provide
no consistent independent support for the image-plane account, and the small
sample does not support a decisive mechanistic conclusion in either direction.
The internally dated plan and frozen criterion are released, but should not be
confused with an independently timestamped public preregistration. The
CourtDyn interventions remain the basis for the bounded behavioural claim.

% ===========================================================================
